% ── §1.7: Διάρθρωση της Διπλωματικής ────────────────────────

\section{Διάρθρωση της Διπλωματικής}
\label{sec:outline}

Το υπόλοιπο της εργασίας διαρθρώνεται ως εξής:

Το \textbf{Κεφάλαιο~\ref{ch:background}} παρουσιάζει το απαραίτητο
θεωρητικό υπόβαθρο: αρχές μηχανικής μάθησης και 1D-CNN για ανάλυση
χρονοσειρών, αρχιτεκτονική και αλγορίθμους Ομοσπονδιακής Μάθησης,
αρχές Αναγνώρισης Δραστηριότητας, έννοια Ψηφιακού Διδύμου, θεωρία
παιγνίων Stackelberg και Shapley Values.

Το \textbf{Κεφάλαιο~\ref{ch:data_analysis}} αναλύει τα σύνολα δεδομένων
που χρησιμοποιούνται στην εργασία (UCI HAR και AAL), παρουσιάζει
Διερευνητική Ανάλυση Δεδομένων (EDA) με έμφαση στην ετερογένεια χρηστών
και περιγράφει τον αγωγό προ-επεξεργασίας (preprocessing pipeline).

Το \textbf{Κεφάλαιο~\ref{ch:implementation}} περιγράφει λεπτομερώς την
αρχιτεκτονική του 1D-CNN, την υλοποίηση του FL πλαισίου εξ ολοκλήρου σε PyTorch,
και τον σχεδιασμό του μηχανισμού κινήτρων Stackelberg.

Το \textbf{Κεφάλαιο~5} παρουσιάζει τις μεθόδους
αξιολόγησης συνεισφοράς clients: μετρικές βάσει όγκου, LOO, Shapley
Values και gradient-based προσέγγιση.

Το \textbf{Κεφάλαιο~6} παρουσιάζει τη ρύθμιση των
πειραμάτων και τα αποτελέσματα: σύγκριση datasets, απόδοση FL,
αποτελέσματα αξιολόγησης συνεισφοράς και επίδραση των κινήτρων.

Το \textbf{Κεφάλαιο~7} συνοψίζει τα συμπεράσματα,
αναλύει ανοικτά ζητήματα ασφάλειας και ιδιωτικότητας, και προτείνει
κατευθύνσεις μελλοντικής έρευνας.
